\section{Conclusions}\label{sec:conclusions}

\subsection{Main Results}\label{sec:main_results_p2}

This paper has established the following results:

\begin{enumerate}
    \item \textbf{Five graph archetypes emerge from the Killing spectrum.} The classification is exact, derived (not postulated), and experimentally verified for $19$ algebras. Each Killing eigenvalue classifies one generator's edge as cyclic ($\Kil_\la < 0$), free ($\Kil_\la > 0$), or decoupled ($\Kil_\la = 0$).

    \item \textbf{Spectral quality: $19/37$ algebras Ramanujan, transition at $h^* \approx 9.5$.} The Ramanujan property provides a sharp criterion for gauge algebra viability. The SM sits deep in the ordered (Ramanujan) phase.

    \item \textbf{Killing as diagnostic, not generative.} RAG experiments demonstrate two-phase crystallization (activation + hardening) and a tempering effect where withdrawal of cycle pressure purifies the Killing signal ($\Kil_{\mathrm{rr}}$ $\AUC$ improves from $0.52$ to $0.98$).

    \item \textbf{Triple balance: aggregation/expansion/entropy as tensegrity.} The three tensions coexist in the Killing spectrum and determine the graph archetype. Formalized as a tensegrity functional $\Tens[\Kil]$.

    \item \textbf{The graph encodes observables through topological invariants.} Five graph-topological building blocks ($\mathcal{R}_2$, $\mathcal{R}_3$, $R_{EE}$, $h_2^{\vee 4}$, $h_3^{\vee 4}$) bridge to the Standard Model.
\end{enumerate}

\subsection{Contributions of Paper 2}\label{sec:contributions}

\begin{itemize}
    \item Dictionary: Killing eigenvalues $\to$ graph topology (derived, not postulated)
    \item Ramanujan property as spectral selection criterion
    \item Ihara zeta of adjoint graphs (first systematic computation)
    \item Experimental evidence for diagnostic vs. generative Killing
    \item Tensegrity functional $\Tens[\Kil]$ proposed
    \item Bridge to Paper 3: five topological building blocks
\end{itemize}

\subsection{Predictions of Paper 2}\label{sec:predictions}

\begin{itemize}
    \item No gauge algebra with dimension $> \sim\!50$ emerges from the variational principle (Ramanujan barrier).
    \item GUT algebras are marginal: su(5) OK, so(10) close, E$_6$ fails.
    \item Ramanujan transition = chaos phase transition ($\langle\mathrm{OTOC}\rangle$ sign change).
    \item Pole migration velocity as order parameter for spectral phase transition.
\end{itemize}

\subsection{Bridge to Paper 3}\label{sec:bridge_p3}

The graph invariants (spectral ratios, Coxeter numbers) are the natural inputs for physical observables. The graph is the computational substrate: the topology generates a discrete lattice in spectral space where the observables live as points. Paper 3 will derive the SM parameters from these topological invariants.

\subsection{The Sequence: Emergence $\to$ Topology $\to$ Observables}\label{sec:sequence}

\begin{itemize}
    \item \textbf{Paper 1} demonstrated that Lie algebras \emph{emerge} (Phase I: Killing metric variational principle).
    \item \textbf{Paper 2} demonstrates that each emergent algebra produces a graph with \emph{specific topology} (Phase II: graph topology + spectral quality).
    \item \textbf{Paper 3} will demonstrate that the graph \emph{computes} the SM parameters (Phase III: Standard Model from graph topology).
\end{itemize}

The sequence is: \textbf{emergence $\to$ topology $\to$ observables}. The five spectral blocks are the bridge: Paper 1 produces them, Paper 2 classifies them spectrally, Paper 3 uses them to predict.
